\documentclass[journal]{IEEEtran}
\begin{document}

\title{Dynamic Kidney Exchange Optimization Under Operational Constraints for Equitable Allocation}

\maketitle
\begin{abstract}
Kidney exchange programs clear their pools repeatedly, yet fairness in
kidney exchange is usually studied within a single clearing round. This
paper examines how the price of fairness evolves over repeated rounds,
whether matching every month or waiting for a larger pool better balances
efficiency and equity, and what a multi-objective evolutionary search adds
to an integer programming benchmark. A synthetic simulator informed by
public statistics of the Australian and New Zealand paired kidney exchange simulates arrivals, withdrawals, exits to deceased-donor transplantation and
transport limits between six regional centers. Each round balances
transplant volume against three fairness criteria and is solved by an
evolutionary algorithm guided by integer programming. Over 25 monthly
rounds and 100 paired runs, the framework transplanted as many patients as
volume maximization while reducing the access gap of highly sensitized
patients by 21\%--24\% and the logistics gap by 28\%--34\%. The price of
fairness averaged under 2\% per round, and most deferred transplants were
made in later rounds. Waiting for a larger pool lost at least 12\% of transplants. The two-stage integer program, solved within practical time
limits, selected fairer maximum-volume matchings, whereas the evolutionary
front offered useful trade-offs below maximum volume. All results are
synthetic.
\end{abstract}

\begin{IEEEkeywords}
Kidney exchange, dynamic matching, price of fairness, evolutionary
multi-objective optimization, adaptive operator selection.
\end{IEEEkeywords}

\section{Introduction}
\label{sec:intro}

\IEEEPARstart{K}{idney} exchange programs (KEPs) enable living-donor
transplants between biologically incompatible donor--patient pairs by
arranging cyclic exchanges and chains initiated by altruistic
donors~\cite{roth2004kidney}.
In August 2019, New Zealand joined the Australian Paired Kidney Exchange
(AKX) to form the trans-Tasman Australian and New Zealand Paired Kidney
Exchange (ANZKX), which involves more than 25 transplant units and
matches pairs continuously through the national OrganMatch system as
new pairs enter the program~\cite{anzkxprotocol2021}.
Between 60 and 90 pairs are registered at any one time, and the
program performs 50--66 transplants per year~\cite{ANZKX2021,DonateLife2024}.

Operational KEPs rely on combinatorial optimization.
Selecting a maximum set of disjoint exchanges is NP-hard when cycles of
length $K \geq 3$ are allowed~\cite{abraham2007clearing}, and integer
linear programs (ILPs), such as the position-indexed chain-edge
formulation (PICEF)~\cite{dickerson2016picef}, are the standard way of
clearing a pool~\cite{biro2021modelling}.
These methods optimize the number of transplants, or a hierarchy of
criteria usually led by it, over one static compatibility graph.

Fairness has mostly been studied in the same static setting.
Dickerson~\textit{et~al.}~\cite{dickerson2014price} defined the price of
fairness (PoF) in kidney exchange as the relative loss of transplants
caused by favouring highly sensitized patients, whose high calculated
panel-reactive antibody (cPRA) levels leave them with few compatible
donors.
Lexicographic and individual notions of fairness
followed~\cite{mcelfresh2018balancing,farnadi2021individual}.
These studies define fairness among patients, and its cost, within a
clearing round.
Real programs, however, clear the pool repeatedly: unmatched patients
keep waiting, some withdraw, some receive a deceased-donor kidney, and
new pairs arrive.

The dynamic side of kidney exchange has been studied mainly in terms of
efficiency~\cite{unver2010dynamic,akbarpour2020thickness,dickerson2015futurematch};
in particular, the value of waiting for a thicker pool depends on
information about departures~\cite{akbarpour2020thickness}. In this paper,
we track fairness and the efficiency given up for it across sequential
clearing rounds.

\begin{figure}[!t]
  \centering
  \includegraphics[width=\columnwidth]{fig1_anzkx.jpg}
  \caption{Dynamic pool model. (a)~Six regional nodes and illustrative
transport routes; all are feasible under the 24-hour cold ischemia limit.
(b)~Monthly arrivals ($\lambda = 10$), deceased-donor exits
($\mu_d = 0.02$) and withdrawal after $\tau \in \{6, 12, 18\}$ months.}
  \label{fig:anzkx}
\end{figure}

This study extends earlier work on fairness-aware optimization for ANZKX that considered a static pool cleared once. Here we move to a dynamic pool: pairs arrive and withdraw, patients
may leave for a deceased-donor transplant, and kidneys must be transported
between distant centers within a cold ischemia time (CIT) limit, which in
ANZKX usually exceeds 10 hours for trans-Tasman exchanges and for
exchanges between the east and west coasts of
Australia~\cite{anzkxprotocol6}.  Fig.~\ref{fig:anzkx} summarizes the
model.
The central research question is: \emph{how does the price of fairness
evolve across sequential dynamic clearing rounds in an ANZKX-informed
simulator under operational constraints, and which clearing
strategy, myopic or batched, best preserves the efficiency-equity
trade-off over time?}
We address six research questions: how the per-round price of fairness evolves over rounds
(RQ1); how myopic and batched clearing compare (RQ2); how the
framework compares with the baselines, including a two-stage ILP that seeks the fairest maximum matching (RQ3); what the
evolutionary front offers below maximum volume (RQ4); what operator
selection and ILP guidance contribute (RQ5); and how sensitive and
scalable the results are (RQ6). All experiments use synthetic pools generated from public ANZKX information; no patient-level data were used.

The main contributions of this paper are as follows.
\begin{enumerate}
  \item \textbf{A price-of-fairness trajectory} over 25 clearing rounds,
    showing that the per-round price is small and largely recovered in
    later rounds.
  \item \textbf{A comparison of myopic and batched clearing} on fairness as
    well as volume, under withdrawals and competing exits.
  \item \textbf{An ANZKX-informed dynamic simulator} with Poisson
    arrivals, virtual crossmatching, withdrawal, a deceased-donor competing
    exit and CIT limits between six regional nodes, using common random
    numbers so that every method faces the same patients.
  \item \textbf{An ILP-guided NSGA-III, assessed against a two-stage ILP
    benchmark} at three levels of comparison, separating the quality of the
    front, the implemented matching and the ILP solutions obtained within
    time limits.
  \item \textbf{A broad evaluation} of ten configurations at three
    withdrawal thresholds with 100 paired runs, including selection-rule,
    sensitivity, scalability, cost and out-of-distribution tests on PrefLib
    dataset~00036~\cite{mattei2013preflib}.
\end{enumerate}
Section~\ref{sec:related} reviews related work, Section~\ref{sec:problem}
formulates the problem, Section~\ref{sec:method} describes the framework,
Sections~\ref{sec:experiments} and~\ref{sec:results} present the
experiments and results, and Section~\ref{sec:conclusion} concludes.

\section{Related Work}
\label{sec:related}

We review four streams of work; Table~\ref{tab:related} summarizes how
representative studies cover the design criteria of our framework.

\subsection{Integer Programming for Kidney Exchange}
\label{sec:related:ilp}

Clearing a kidney exchange pool is NP-hard when cycles of length
$K \geq 3$ are allowed~\cite{abraham2007clearing}. Three-way exchanges
add substantially to the number of transplants, while longer ones add
little~\cite{roth2007efficient,saidman2006increasing}, and non-directed
donors can start chains~\cite{montgomery2006domino}. Formulations have
since become more compact~\cite{constantino2013new} and better at
handling long chains and multiple
criteria~\cite{glorie2014kidney,anderson2015new}, including the
position-indexed PICEF model~\cite{dickerson2016picef}, branch-and-price
with decision diagrams~\cite{riascos2024branch} and
branch-and-cut-and-price~\cite{lam2020branch}; two surveys compare
them~\cite{barkel2025operational,biro2021modelling}. ANZKX matches pairs
through the national OrganMatch system~\cite{anzkxprotocol2021}. These
methods optimize the number of transplants, or a hierarchy of criteria
led by it~\cite{biro2021modelling}, within a single round; we instead
treat fairness and efficiency as separate objectives.

\subsection{Dynamic Kidney Exchange}
\label{sec:related:dynamic}

ANZKX matches pairs continuously as they register~\cite{ANZKX2021}.
\"{U}nver~\cite{unver2010dynamic} characterized efficient dynamic
mechanisms, Dickerson~\textit{et~al.}~\cite{dickerson2012optimizing}
found that a chain cap of four performed best in dynamic experiments, which motivates our $L = 4$, and FutureMatch~\cite{dickerson2015futurematch}
learned weights for dynamic matching. Glorie~\textit{et~al.}~\cite{glorie2022health}
compared allocation policies by quality-adjusted life-years in a
simulation of the Dutch program. On when to clear a pool,
greedy matching is near optimal for waiting
time~\cite{anderson2017efficient}, waiting pays off only when departures
can be anticipated~\cite{akbarpour2020thickness}, and related work
studies unpaired exchange~\cite{akbarpour2025unpaired} and online
stochastic matching~\cite{huang2021online}. Across hospitals and
countries, studies address participation
incentives~\cite{ashlagi2014free,hajaj2015strategy} and fairness among
programs over time~\cite{klimentova2021fairness,benedek2025partitioned,druzsin2024enckep}.
We instead optimize several fairness criteria for patients across
consecutive rounds.

\subsection{Fairness in Kidney Exchange}
\label{sec:related:fairness}

The price of fairness measures the efficiency lost when a fair
allocation replaces an efficient
one~\cite{bertsimas2011price,caragiannis2012efficiency}.
Dickerson~\textit{et~al.}~\cite{dickerson2014price} brought it to kidney exchange for highly sensitized patients, showing that it is small in the standard theoretical model and quantifying the trade-off on real exchange data. Later work bounded the efficiency loss
of lexicographic priority~\cite{mcelfresh2018balancing}, derived patient
weights from human judgments~\cite{freedman2020adapting}, used multiple
optimal solutions or lotteries for individual
fairness~\cite{farnadi2021individual,starnaud2025adaptation,chang2026optimizing},
and required outcomes independent of protected
features~\cite{zhang2026fairness}; European programs also differ in their
criteria~\cite{biro2021modelling}. These studies measure fairness within
a single round. We treat waiting-time equity, access for highly
sensitized patients and logistics disparity as separate objectives and
follow them across rounds.

\begin{table*}[!htbp]
\caption{Comparison of representative kidney exchange studies against
  key design criteria. \checkmark~= criterion explicitly addressed;
  {--}~= absent or not explicit.
  Groups: (A)~static ILP, (B)~dynamic KEP, (C)~fairness-aware KEP.}
\label{tab:related}
\begin{tabular}{llcclcc}
\toprule
\textbf{Work} & \textbf{Year} & \textbf{Dyn.\ Pool} &
  \textbf{Obj.} & \textbf{Fairness Dimension} &
  \textbf{AOS} & \textbf{hi-cPRA} \\
\midrule
\GROUP{(A) Static ILP Formulations} \\
Abraham~\textit{et~al.}~\cite{abraham2007clearing}
  & 2007 & --          & 1 & None                        & -- & -- \\
Constantino~\textit{et~al.}~\cite{constantino2013new}
  & 2013 & --          & 1 & None                        & -- & -- \\
Anderson~\textit{et~al.}~\cite{anderson2015new}
  & 2015 & --          & 1 & None                        & -- & -- \\
Dickerson~\textit{et~al.}~\cite{dickerson2016picef}
  & 2016 & --          & 1 & None                        & -- & -- \\
Riascos-\'{A}lvarez~\textit{et~al.}~\cite{riascos2024branch}
  & 2024 & --          & 1 & None                        & -- & -- \\
\GROUP{(B) Dynamic KEP} \\
\"{U}nver~\cite{unver2010dynamic}
  & 2010 & \checkmark  & 1 & None                        & -- & -- \\
Dickerson~\textit{et~al.}~\cite{dickerson2012optimizing}
  & 2012 & \checkmark  & 1 & None                        & -- & -- \\
Dickerson and Sandholm~\cite{dickerson2015futurematch}
  & 2015 & \checkmark  & 1 & Human value judgements      & -- & -- \\
Anderson~\textit{et~al.}~\cite{anderson2017efficient}
  & 2017 & \checkmark  & 1 & None                        & -- & -- \\
Akbarpour~\textit{et~al.}~\cite{akbarpour2020thickness}
  & 2020 & \checkmark  & 1 & None                        & -- & -- \\
\GROUP{(C) Fairness-Aware KEP} \\
Dickerson~\textit{et~al.}~\cite{dickerson2014price}
  & 2014 & --          & 1 & Lexicographic, weighted     & -- & \checkmark \\
McElfresh and Dickerson~\cite{mcelfresh2018balancing}
  & 2018 & --          & 2 & Hybrid lexicographic        & -- & \checkmark \\
Freedman~\textit{et~al.}~\cite{freedman2020adapting}
  & 2020 & --          & 1 & Elicited human values       & -- & -- \\
Farnadi~\textit{et~al.}~\cite{farnadi2021individual}
  & 2021 & --          & 1 & Individual                  & -- & -- \\
Zhang~\textit{et~al.}~\cite{zhang2026fairness}
  & 2026 & --          & 1 & Calibration-based           & -- & -- \\
\textbf{This work}
  & \textbf{2026} & \checkmark & \textbf{4}
  & \textbf{Waiting-time Gini, hi-cPRA gap, logistics}
  & \checkmark & \checkmark \\
\bottomrule
\end{tabular}
\end{table*}

\subsection{Multi-Objective Evolutionary Algorithms and Adaptive
Operator Selection}
\label{sec:related:moea}

NSGA-II~\cite{deb2002nsga2} and NSGA-III~\cite{deb2014nsga3} are standard
Pareto-based algorithms, the latter using reference points for four or
more objectives; MOEA/D~\cite{zhang2007moead,trivedi2017survey} is the
main decomposition-based alternative, and the hypervolume reference point
affects comparisons~\cite{ishibuchi2017reference}. Adaptive operator
selection chooses operators or parameter values during the
search~\cite{thierens2005adaptive}, often with bandits built on
UCB1~\cite{auer2002finite,maturana2009extreme,fialho2010analyzing,li2014frrmab}.
Methods are stateless or state-based~\cite{pei2025aos}; deep Q-networks
drive recent state-based
selection~\cite{yin2024dueling,ming2024constrained}, and contextual
bandits such as LinUCB~\cite{li2010contextual} use a context vector
without a neural network. Deep reinforcement learning needs training
before use~\cite{li2021deeprl}, whereas bandits adapt online, which suits
a pool that changes every month; fairness has also been linked to
multi-objective optimization more broadly~\cite{yu2025fairness}. We guide
NSGA-III with the volume-optimal ILP matching, compare stateless and
contextual bandits for mutation-rate selection, and assess the result
against a two-stage ILP benchmark.

\section{Problem Formulation}
\label{sec:problem}

\subsection{Dynamic Pool Model}
\label{sec:problem:model}

We model the exchange in monthly rounds $t = 1,\dots,T$. The pool $V_t$
contains incompatible donor--patient pairs and altruistic donors, and
$P_t \subseteq V_t$ is the set of patients. Each patient $i$ has a blood
type, a cPRA value, a waiting time $w_i$ in months, and a home center
$h(i) \in \{\text{Perth, Adelaide, Melbourne, Sydney, Brisbane, Auckland}\}$. Highly sensitized patients form the set
$S_t = \{ i \in P_t : \cpra_i \geq 0.80 \}$.

After the matching $M_t$ is implemented, the pool evolves as
\begin{equation}
  V_{t+1} = \bigl(V_t \setminus M_t \setminus X_t \setminus W_t\bigr)
            \cup A_t ,
  \label{eq:pool}
\end{equation}
applied in this order. Matched pairs leave. Each remaining patient
leaves for a deceased-donor transplant with probability $\mu_d$ (the
competing exit $X_t$), since patients in the Australian paired donation
program are often also listed for a deceased-donor
kidney~\cite{Ferrari2012}. Waiting times then grow by one month, pairs
that have waited $\tau$ months withdraw ($W_t$), and
$|A_t| \sim \mathrm{Poisson}(\lambda)$ new pairs arrive. The threshold
$\tau$ is a modeling assumption informed by~\cite{Ferrari2012}, and we
vary it over $\{6, 12, 18\}$ months.

New pairs follow the Saidman model~\cite{saidman2006increasing} with
three changes. Blood types follow Australian population frequencies.
Thirty percent of patients have $\cpra \in [0.80, 0.99]$, and the rest
are split equally between $[0, 0.05)$ and $[0.05, 0.80)$. A pair joins
only if its donor cannot give to its own patient, because of ABO
incompatibility or a positive crossmatch. Eight percent of arrivals are
altruistic donors.

\subsection{Compatibility and Logistics}
\label{sec:problem:compat}

The compatibility graph $G_t = (V_t, E_t)$ has an arc $(i, j)$ when the
donor of $i$ can give to the patient of $j$:
\begin{equation}
  (i,j) \in E_t \iff \mathrm{ABO}(i,j) \wedge \mathrm{XM}^{-}(i,j)
  \wedge T^{\mathrm{tr}}_{h(i)h(j)} + T^{\mathrm{prep}} \leq \cit_{\max} .
\label{eq:edge}
\end{equation}
The virtual crossmatch $\mathrm{XM}^{-}(i,j)$ is negative with
probability $1 - \cpra_j$. It is drawn once per donor--patient couple,
so the result is the same in every round and configuration. The
transport time $T^{\mathrm{tr}}_{ab}$ between two centers is the flight
time plus one hour of ground transport at each end and a wait for a
flight (2~h domestic, 4~h trans-Tasman); within a center it is zero. We
set $T^{\mathrm{prep}} = 6$~h and $\cit_{\max} = 24$~h. These transport
and CIT values are modeling assumptions, not ANZKX operational
parameters. HLA profiles are random, and each arc carries a weight
$\omega_{ij} = 1 - m_{ij}/6$, where $m_{ij}$ is the number of HLA
mismatches; it is used only to break ties.

\begin{figure}[!t]
\centering
\includegraphics[width=\columnwidth]{fig2_structures.jpg}
\caption{Exchange structures: (a)~two-way and (b)~three-way cycles, and
(c)~a chain started by an altruistic donor (A) with up to four
transplants. Arrows point from donor (D) to recipient (P).}
\label{fig:structures}
\end{figure}

A matching consists of vertex-disjoint cycles of at most $K = 3$ pairs
and chains of at most $L = 4$ transplants, each started by an altruistic
donor (Fig.~\ref{fig:structures}). All cycles and chains are enumerated
once per round, and every method uses the same list. We keep at most
5{,}000 chains per round: all chains with one or two transplants and a
random sample of longer ones. When this cap binds, we add the chains of
the exact optimum found with PICEF~\cite{dickerson2016picef}, so the
list always contains a volume-optimal matching.

\subsection{Objective Functions}
\label{sec:problem:objectives}

For a matching $M$, let $P = P_t$ and $S = S_t$, and let $R$ and $U$ be
the patients at the distant centers (Perth, Auckland) and at the central
ones (Adelaide, Melbourne, Sydney, Brisbane). Altruistic donors are not
counted in any of these sets.

\begin{definition}[Transplant volume]
$F_1(M) = |M \cap P|$, the number of patients transplanted.
\end{definition}

\begin{definition}[Waiting-time equity]
$F_2(M)$ is the Gini coefficient of the waiting times after the round,
where matched patients keep $w_i$ and the others move to
$w'_i = w_i + 1$:
\begin{equation}
  F_2(M) = \frac{\sum_{i \in P}\sum_{j \in P} |w'_i - w'_j|}
                {2 |P|^2 \, \overline{w'}} ,
  \label{eq:gini}
\end{equation}
where $\overline{w'}$ is the mean of $w'$, and $F_2 = 0$ when $|P| < 2$
or all waiting times are zero.
\end{definition}

\begin{definition}[High-cPRA access gap]
\begin{equation}
  F_3(M) = \left| \frac{|M \cap S|}{|S|} - \frac{|M \cap P|}{|P|}
  \right| ,
  \label{eq:f3}
\end{equation}
with $F_3 = 0$ when $S = \emptyset$.
\end{definition}

\begin{definition}[Logistics disparity]
\begin{equation}
  F_4(M) = \left| \frac{|M \cap R|}{|R|} - \frac{|M \cap U|}{|U|}
  \right| ,
  \label{eq:f4}
\end{equation}
with $F_4 = 0$ when $R$ or $U$ is empty.
\end{definition}

$F_3$ compares the match rate of highly sensitized patients with the
overall rate, and the absolute value penalizes both neglect and
over-correction. A smaller $F_3$ therefore means more equal rates, not
necessarily more transplants for these patients. $F_4$ is a modeled
proxy for geographic disparity, not a validated measure of inequity.
Each round solves
\begin{equation}
  \min_{M \in \mathcal{M}_t} \;
  \bigl( 1 - F_1/|P|,\; F_2,\; F_3,\; F_4 \bigr) ,
  \label{eq:mop}
\end{equation}
where $\mathcal{M}_t$ is the set of feasible matchings and all four
objectives lie in $[0, 1]$. The empty matching has $F_3 = F_4 = 0$ and
is Pareto-optimal, so it is never implemented.

\subsection{Price of Fairness}
\label{sec:problem:pof}

\begin{definition}[Price of fairness]
Let $F_1^{\ilp}(t)$ be the largest number of transplants possible in
round $t$ and $F_1(t)$ the number in the implemented matching. Then
\begin{equation}
  \pof = \frac{F_1^{\ilp}(t) - F_1(t)}{F_1^{\ilp}(t)} \times 100\% ,
  \label{eq:pof}
\end{equation}
which is undefined when $F_1^{\ilp}(t) = 0$.
\end{definition}

We summarize $\{\pof\}_{t=1}^{T}$ by its mean over the rounds where it
is defined, its variance, and its per-run slope. Because $\pof$
compares two matchings of the same pool, and a patient passed over in
one round may be matched in the next, we also measure the loss over the
whole trajectory:
\begin{equation}
  \Delta = 1 - \frac{\sum_t F_1(t)}{\sum_t F_1^{\mathrm{MI}}(t)} ,
  \label{eq:delta}
\end{equation}
where $F_1^{\mathrm{MI}}(t)$ is the number of transplants made by the
myopic ILP in the run with the same seed. Unlike $\pof$, $\Delta$ also
compares clearing policies, since batched clearing skips rounds.

\subsection{Integer Programming Baselines}
\label{sec:problem:ilp}

Let $x_k \in \{0,1\}$ select structure $k$, which has $|k|$
transplants. The volume-maximizing ILP is
\begin{equation}
  \max \sum_k \bigl(|k| + \epsilon\, \bar\omega_k\bigr) x_k
  \quad \text{s.t.} \quad
  \sum_{k \ni v} x_k \leq 1 \;\; \forall v \in V_t ,
  \label{eq:ilp}
\end{equation}
where $\bar\omega_k$ is the normalized HLA weight of structure $k$ and
$\epsilon = 10^{-3}$ breaks ties without losing a transplant. The
lexicographic baseline of
Dickerson~\textit{et~al.}~\cite{dickerson2014price} first maximizes the
number of highly sensitized patients transplanted, then the total
number of transplants.

The \emph{two-stage ILP} formulation seeks the fairest matching among
those of maximum volume. Stage~1 is~\eqref{eq:ilp}, which gives $F_1^{*}$ and
the matching $M^{\ilp}_t$. Stage~2 keeps the same constraints, fixes
$\sum_k |k|\, x_k = F_1^{*}$, and minimizes
\begin{equation}
  \frac{F_2}{r_2} + \frac{F_3}{r_3} + \frac{F_4}{r_4}
  - \epsilon' \sum_k \bar\omega_k x_k ,
  \quad r_k = \max\bigl(F_k(M^{\ilp}_t),\, 0.01\bigr) ,
  \label{eq:twostage}
\end{equation}
so each objective is scaled by its value in the volume-maximizing
matching, and $\epsilon' = 10^{-7}$ breaks ties by HLA quality. With the
volume fixed, the overall match rate and $\sum_i w'_i$ are constants, so
$F_3$, $F_4$ and the Gini coefficient $F_2$ become linear after adding
auxiliary variables. The formulation is therefore exact: its optimal solution is
the fairest maximum matching under~\eqref{eq:twostage}. In practice every
solve has a time limit (Table~\ref{tab:params}). When the limit is reached,
the solver returns the best solution found, which is not proven optimal; we
therefore report how often each experiment reached its limit and refer to
such results as best-found solutions. A second variant minimizes
$F_3$ alone, with $F_4$ and HLA quality as tie-breakers. Stage~2 can
also be solved with the volume fixed below $F_1^{*}$, which we use to
measure the fairness available at lower volumes. All ILPs are solved
with CBC, with HiGHS as a fallback.

\section{Proposed Framework}
\label{sec:method}

\subsection{Overview}
\label{sec:method:overview}

Fig.~\ref{fig:framework} shows one round of the framework, and Algorithm~\ref{alg:main} gives a full trajectory. A clearing policy
decides \emph{when} the pool is cleared, and an ILP-guided NSGA-III
decides \emph{how}. In each round that is cleared, the ILP~\eqref{eq:ilp}
first gives the maximum number of transplants $F_1^{\ilp}(t)$. NSGA-III
then approximates the Pareto front of~\eqref{eq:mop}, a selection rule
implements one matching from it, and the pool evolves
by~\eqref{eq:pool}. Inside the search, a bandit sets the mutation rate
at every generation. The front is the main output of a round: it holds
matchings at maximum volume and below it, from which a selection rule
chooses according to how much volume may be given up for fairness.

\begin{figure*}[!htbp]
\centering
\includegraphics[width=\textwidth]{fig3_framework.png}
\caption{One round of the framework. The ILP solution guides NSGA-III, a
bandit sets the mutation rate, and small rounds are solved exactly.}
\label{fig:framework}
\end{figure*}

\ALGORITHM

\subsection{Encoding, Repair and Variation}
\label{sec:method:nsga3}

A solution is a binary vector with one bit per vertex, marking whether
the pair or altruistic donor is eligible to be matched. A structure in
$\mathcal{K}_t$ is usable only if all its vertices are eligible. The
repair operator scans the usable structures in a fixed priority order:
structures of the ILP solution first (Section~\ref{sec:method:hybrid}),
then structures with more transplants, then those with higher HLA
quality. It accepts each structure that is disjoint from those already
accepted, and writes the result back into the individual (Lamarckian
repair). Every individual therefore encodes a feasible matching over
the same structure list as the ILP.

NSGA-III~\cite{deb2014nsga3} evolves $N = 100$ individuals for
$G_{\max} = 150$ generations per round, with Das--Dennis reference
directions ($p = 12$ divisions, 455 directions for four objectives).
Offspring are produced by uniform crossover with probability
$p_c = 0.9$ and bit-flip mutation with rate $p_m$, chosen by the bandit.
Duplicates are removed.

\subsection{ILP Guidance, Archive and Exact Small Rounds}
\label{sec:method:hybrid}

Packing cycles and chains is NP-hard~\cite{abraham2007clearing}, and
NSGA-III alone rarely reaches the largest matching. In 30 pools of 75
vertices it found 32.6 transplants on average against an ILP optimum of
42.3, and never reached the optimum (Section~\ref{sec:results:rq5}). Because
$\pof$ is measured against the ILP optimum, such search failure would
inflate it. We therefore insert the ILP solution into the initial
population and give its structures priority in the repair. An external
archive collects every non-dominated solution of the whole population
at every generation, so niching cannot discard the volume-optimal
matching. The archive forms the final front and always contains the ILP
optimum. This reduces the bias that search failure adds to $\pof$, but
part of the remaining loss is still set by the selection rule
(Section~\ref{sec:method:select}). All six evolutionary configurations
use the same guidance. It is the component that makes $\pof$ measurable.

In dynamic configurations, the previous round's front is also inserted
into the initial population, after removing the vertices that have left
the pool. The ILP solution and these warm-start solutions fill at most
half of the population; each remaining individual is random, with its
own eligibility rate drawn uniformly from $[0.2, 1]$.

After matches, withdrawals and exits, a round may contain few
structures. If $\mathcal{K}_t$ admits at most $2N = 200$ matchings,
including the empty one, all are enumerated and the exact Pareto front
is used. The bandit is not called in such rounds, since the population
could not be filled with distinct individuals; its state carries over
to the next round.

\subsection{Adaptive Operator Selection}
\label{sec:method:aos}

The bandit chooses among the mutation rates
$\mathcal{A} = \{0.01, 0.05, 0.10, 0.25\}$. After generation $g$, the
chosen arm receives
\begin{equation}
  r_g = \frac{\Delta\mathrm{HV}_g}{\max_{g' \leq g} \Delta\mathrm{HV}_{g'}}
        \in [0, 1], \quad
  \Delta\mathrm{HV}_g = \max\bigl(0,\, \mathrm{HV}_g - \mathrm{HV}_{g-1}\bigr),
\label{eq:reward}
\end{equation}
with $r_g = 0$ until a gain is observed in the round. $\mathrm{HV}_g$ is
the hypervolume of the non-dominated individuals in the population,
computed on~\eqref{eq:mop} with reference point $(1.1, 1.1, 1.1, 1.1)$.
The reward is thus the hypervolume gain of a generation relative to the
largest gain so far in the round.

\subsubsection{Stateless UCB1}
UCB1~\cite{auer2002finite} considers the last $\ell = 15$ choices and
selects
\begin{equation}
  a_g = \arg\max_{a \in \mathcal{A}} \;
        \hat{Q}_a + \kappa \sqrt{\frac{2 \ln n}{n_a}} ,
  \label{eq:ucb1}
\end{equation}
where $\hat{Q}_a$ and $n_a$ are the mean reward and number of plays of
arm $a$ in the window, $n = \sum_a n_a$, and $\kappa = 1$. An arm absent
from the window is played first.

\subsubsection{Contextual LinUCB}
LinUCB~\cite{li2010contextual} also observes the pool state
\begin{equation}
  \mathbf{x}_t = \left[\frac{|P_t|}{150},\;
  \frac{|S_t|}{|P_t|},\;
  \frac{|X_{t-1}|}{|P_t|},\;
  \frac{t}{T},\;
  \mathrm{Gini}_t,\;
  \frac{\mathrm{PoF}_{t-1}}{100}\right]^{\top}  ,
  \label{eq:context}
\end{equation}
with each feature clipped to $[0,1]$, namely the pool size, the share of highly sensitized patients, the
share of competing exits in the previous round, the progress through
the horizon, the waiting-time Gini coefficient of the pool, and the
previous price of fairness. It selects
\begin{equation}
  a_g = \arg\max_{a \in \mathcal{A}} \;
        \boldsymbol{\theta}_a^{\top} \mathbf{x}_t
        + \alpha \sqrt{\mathbf{x}_t^{\top} \mathbf{D}_a^{-1} \mathbf{x}_t} ,
  \label{eq:linucb}
\end{equation}
with $\alpha = 1$, where $\mathbf{D}_a$ and $\boldsymbol{\theta}_a$ are
updated by ridge regression on the observed rewards. Since
$\mathbf{x}_t$ is fixed within a round, LinUCB adapts across rounds,
whereas UCB1 responds to rewards alone. Dynamic configurations keep the
bandit's state across rounds; static configurations create a new bandit
in every round and use no warm start.

\subsection{Implemented Solution}
\label{sec:method:select}

The empty matching is removed from the final front, and only solutions
with $F_1 \geq 0.9 \max F_1$ are kept. Among them, the one closest to
the ideal point in $(F_2, F_3, F_4)$, after min--max scaling over the
kept set, is implemented. Since the front contains the ILP optimum,
$\max F_1 = F_1^{\ilp}(t)$ and
\begin{equation}
  0 \leq \pof \leq 10\% \quad \text{in every cleared round.}
  \label{eq:pofbound}
\end{equation}
When $F_1^{\ilp}(t) < 10$, giving up one transplant would exceed the
10\% allowance, so $\pof = 0$ and the rule picks a volume-optimal
matching. This rule is not the stage-2 objective~\eqref{eq:twostage} of
the two-stage ILP, so at maximum volume the two can choose different
matchings. The quality of the front and the quality of the implemented
matching are therefore distinct, and we report both. The floor acts as a decision parameter: it sets how much
volume may be traded for fairness. The experiments therefore also
evaluate a 100\% floor (volume-optimal matchings only), an 80\% floor,
and the compromise solution over all non-empty front solutions.

\subsection{Clearing Policies}
\label{sec:method:clearing}

\emph{Myopic} clearing ($\pi_M$) clears in every round. \emph{Batched}
clearing ($\pi_B$) clears only when at least $n^{*} = 90$ patients are
waiting, the upper end of the ANZKX pool size~\cite{ANZKX2021}. When the
pool cannot reach $n^{*}$, as at $\tau = 6$, $\pi_B$ never clears. We
therefore also evaluate a variant that clears at least every three
rounds, and one with $n^{*} = 60$, the lower end of the ANZKX pool size.

\section{Experimental Setup}
\label{sec:experiments}

\subsection{Simulation Parameters}
\label{sec:exp:pool}

Table~\ref{tab:params} lists the parameters. We calibrated the arrival
rate by simulating the myopic ILP. With $\lambda = 10$ it performed
64.9--70.0 transplants per year (ANZKX: 50--66), and the pool settled at
31.5, 50.3 and 64.0 patients for $\tau = 6$, 12 and 18 (ANZKX: 60--90).
No single rate matches both targets, so we also test $\lambda = 7$,
which matches the volume, and $\lambda = 13$, which matches the pool
size. With $\cit_{\max} = 24$~h every center pair is feasible; the limit
first binds at 16~h (Perth--Auckland) and rules out 73\% of center
pairs at 12~h. Only the high-cPRA share and the ANZKX volume and pool
size come from published sources. The exit rate, withdrawal thresholds,
transport and CIT constraints and HLA model are modeling assumptions,
not validated ANZKX parameters, and the most influential of them are
varied in the sensitivity analysis.

\begin{table}[!htbp]
\caption{Simulation and algorithm parameters (sensitivity values in
parentheses). Basis: P = published evidence, C = calibrated to
published aggregate statistics, A = modeling assumption or design
choice.}
\label{tab:params}
\begin{tabular}{llc}
\toprule
Parameter & Value & Basis \\
\midrule
Rounds $T$ / round length & 25 / one month & A \\
Initial pool & 75 vertices, waits $\sim U\{0,\dots,5\}$ & A \\
Arrival rate $\lambda$ & 10 per month (7, 13) & C~\cite{ANZKX2021,DonateLife2024} \\
Withdrawal threshold $\tau$ & 6, 12, 18 months & A~\cite{Ferrari2012} \\
Competing exit rate $\mu_d$ & 0.02 per month (0.01, 0.04) & A~\cite{DonateLife2024} \\
High-cPRA share & 30\% of patients & P~\cite{ANZDATA2026} \\
Altruistic donors & 8\% of arrivals & A \\
Cycle / chain limit & $K = 3$ / $L = 4$ & A \\
$\cit_{\max}$ / $T^{\mathrm{prep}}$, transport & 24~h (16, 12) / 6~h, flight-based & A \\
HLA compatibility & random profiles, tie-break only & A \\
Batching threshold $n^*$ & 90 patients (60) & A~\cite{ANZKX2021} \\
NSGA-III & $N = 100$, $G_{\max} = 150$, $p = 12$ & A \\
Operators & uniform crossover $p_c = 0.9$, bit-flip & A \\
Mutation arms & $\{0.01, 0.05, 0.10, 0.25\}$ & A \\
UCB1 / LinUCB & $\ell = 15$, $\kappa = 1$ / $\alpha = 1$, $d = 6$ & A \\
Selection floor & 90\% (100\%, 80\%, compromise) & A \\
Two-stage ILP time limit & 60~s per round (300~s identical pools; 120~s lower levels) & A \\
Runs & 100 paired (30 for sensitivity) & A \\
\bottomrule
\end{tabular}
\end{table}

\subsection{Baselines}
\label{sec:exp:baselines}

Table~\ref{tab:configs} lists the ten configurations. Configurations 1
and 2 solve each round independently, as in a static framework. Configurations 3--6 are the proposed framework under myopic
and batched clearing. Configurations 7 and 8 are the volume-maximizing
ILP and the lexicographic rule of
Dickerson~\textit{et~al.}~\cite{dickerson2014price}, and 9 and 10 are
the two-stage ILPs of Section~\ref{sec:problem:ilp}. All methods use
the same structure list in every round.

\begin{table}[!htbp]
\caption{Compared configurations. All MOEA configurations use ILP
guidance.}
\label{tab:configs}
\begin{tabular}{clllc}
\toprule
\# & Short name & Method & Mutation rate & Clearing \\
\midrule
1 & S-NSGA-II + UCB1  & Static NSGA-II   & UCB1         & Myopic \\
2 & S-NSGA-III fixed  & Static NSGA-III  & $p_m = 0.10$ & Myopic \\
3 & D-UCB1 myopic     & Dynamic NSGA-III & UCB1         & Myopic \\
4 & D-LinUCB myopic   & Dynamic NSGA-III & LinUCB       & Myopic \\
5 & D-UCB1 batched    & Dynamic NSGA-III & UCB1         & Batched \\
6 & D-LinUCB batched  & Dynamic NSGA-III & LinUCB       & Batched \\
7 & Myopic ILP        & ILP~\eqref{eq:ilp} & --         & Myopic \\
8 & Lex. rule & Lexicographic ILP~\cite{dickerson2014price} & -- & Myopic \\
9 & Two-stage (sum)   & Two-stage ILP~\eqref{eq:twostage} & --  & Myopic \\
10 & Two-stage ($F_3$) & Two-stage ILP, $F_3$ only & --      & Myopic \\
\bottomrule
\end{tabular}
\NOTE{Static: fresh bandit and no warm start in every round.}
\NOTE{Dynamic: bandit state and warm start carried across rounds.}
\end{table}

\subsection{Metrics and Statistical Analysis}
\label{sec:exp:metrics}

For each run we record the mean, variance and slope of $\pof$, the
trajectory loss $\Delta$~\eqref{eq:delta}, total and high-cPRA
transplants, high-cPRA withdrawals, $F_2$--$F_4$, hypervolume, the arms
selected, and the time per cleared round. Run $r$ uses the same seed in
every configuration, so all comparisons are paired. A Friedman test is
followed by Wilcoxon signed-rank tests with Bonferroni correction, and
effect sizes are reported as Cohen's $d_z$. The slope of $\pof$ is
tested against zero over all rounds and over rounds 2--25.

\subsection{Comparison with the Two-Stage ILP}
\label{sec:exp:levels}

We compare the framework with the two-stage ILP at three levels, each
answering a different question:
\begin{itemize}
  \item \emph{Identical pools:} the 30 initial pools of 75 vertices are
    cleared once by every method. This tests how close the
    framework's maximum-volume matching is to the two-stage ILP solution on
    the same instance.
  \item \emph{Fixed trajectories:} in every round of the framework's
    own trajectories (UCB1, myopic, 30 runs), the two-stage ILP is also
    solved on the same pool, without changing the trajectory. This
    compares the methods on the pools that arise over time.
  \item \emph{Full trajectories:} each method clears its own pool
    sequence over 25 rounds (100 paired runs). This measures the
    combined effect, including how each method shapes future pools.
\end{itemize}
To assess what the front offers below maximum volume, we use levels
$\rho \in \{1.0, 0.95, 0.90, 0.85, 0.80\}$, where level $\rho$ allows
$V = \lceil \rho F_1^{*} \rceil$ or more transplants. For each level we
record the number of front solutions with $F_1 \geq V$, the fairest of
them, the solution the floor rule would implement, and an
ILP reference: the best stage-2 solution found over the volumes from $V$ to
$F_1^{*}$ (time limit 120~s on identical pools and 60~s on trajectory
rounds, all three $\tau$, 30 runs). Solves that reach the limit give
best-found, not proven optimal, references. In rounds solved by enumeration the
front is exact, so we also count how often several maximum matchings
share the optimal stage-2 score, which shows how often the tie-break
decides. Finally, the 100\%, 90\% and 80\% floors and the compromise
rule are run over full trajectories at all three $\tau$ (30 runs).

\subsection{Sensitivity, Scalability and PrefLib}
\label{sec:exp:ood}

Sensitivity analyses use 30 runs at $\tau = 12$ and vary the values in
parentheses in Table~\ref{tab:params}; $n^* = 60$ is also tested at
$\tau = 18$. They also test batched clearing with a three-round maximum
interval at all $\tau$. For scalability, we run 10 rounds and 10 runs
with initial pools of 200 vertices (framework, volume ILP and
lexicographic rule) and 500 vertices (both ILPs only). Two ablations
use the 30 pools of 75 vertices: one compares both bandits with fixed
mutation rates, with and without ILP guidance; the other compares
NSGA-III alone, the ILP solution as a seed only, the seed with repair
priority, and the full hybrid with the archive (UCB1, exact enumeration
disabled).

For out-of-distribution validation we use PrefLib
dataset~00036~\cite{mattei2013preflib}. Half of each instance's
vertices form the initial pool, the rest arrive over the $T$ rounds,
and $n^*$ is scaled to $90/75$ of the initial pool. All configurations, including both two-stage ILPs, are run on 12 instances
of 128--147 vertices with 10 runs each at $\tau = 12$.

\section{Results and Discussion}
\label{sec:results}

Table~\ref{tab:main} summarizes the main experiment. All comparisons
are paired Wilcoxon tests over 100 runs with Bonferroni correction
unless stated otherwise. In 48\% of cleared rounds the feasible set
was small enough to enumerate exactly, the chain cap bound in 9.6\%,
and no implemented matching fell below the 90\% floor.

\begin{table*}[!htbp]
\caption{Trajectory results over 25 rounds (mean $\pm$ sd over 100
paired runs). $\Delta$: transplants lost against the myopic ILP.
High-cPRA lost: highly sensitized patients who withdrew.
D-/S-: dynamic/static NSGA. S-NSGA-III fixed and D-LinUCB batched are
omitted; they differed from S-NSGA-II + UCB1 and D-UCB1 batched by at
most 1.5\% in transplants and 0.01 in $F_2$--$F_4$.}
\label{tab:main}
\begin{tabular}{clcccccccc}
\toprule
$\tau$ & Configuration & PoF (\%) & $\Delta$ (\%) & Transplants
 & High-cPRA tx & $F_2$ & $F_3$ & $F_4$ & High-cPRA lost \\
\midrule
6 & S-NSGA-II + UCB1   & 1.63 $\pm$ 0.68 & $-$0.03 $\pm$ 1.74 & 173.6 $\pm$ 15.5 & 49.5 $\pm$ 6.9 & 0.354 & 0.060 & 0.097 & 25.6 $\pm$ 4.3 \\
 & D-UCB1 myopic      & 1.60 $\pm$ 0.61 & $-$0.06 $\pm$ 1.55 & 173.7 $\pm$ 15.1 & 49.5 $\pm$ 6.9 & 0.355 & 0.060 & 0.098 & 25.5 $\pm$ 4.4 \\
 & D-LinUCB myopic    & 1.59 $\pm$ 0.63 & $-$0.05 $\pm$ 1.74 & 173.7 $\pm$ 15.4 & 49.4 $\pm$ 6.9 & 0.355 & 0.060 & 0.097 & 25.7 $\pm$ 4.3 \\
 & D-UCB1 batched     & --              & 100.0              & 0.0              & 0.0            & --    & --    & --    & 65.6 $\pm$ 7.8 \\
 & Myopic ILP         & 0.00            & 0.00               & 173.6 $\pm$ 15.4 & 48.8 $\pm$ 7.1 & 0.369 & 0.077 & 0.137 & 26.3 $\pm$ 4.7 \\
 & Lexicographic rule & 1.76 $\pm$ 1.27 & 0.15 $\pm$ 1.49    & 173.3 $\pm$ 15.5 & 59.2 $\pm$ 8.0 & 0.368 & 0.095 & 0.135 & 17.5 $\pm$ 4.5 \\
 & Two-stage (sum)    & 0.00            & $-$0.82 $\pm$ 1.48 & 175.0 $\pm$ 15.6 & 51.2 $\pm$ 7.5 & 0.356 & 0.058 & 0.087 & 24.2 $\pm$ 4.2 \\
 & Two-stage ($F_3$)  & 0.00            & $-$0.14 $\pm$ 1.49 & 173.8 $\pm$ 15.3 & 51.4 $\pm$ 7.3 & 0.369 & 0.057 & 0.096 & 24.0 $\pm$ 4.2 \\
\midrule
12 & S-NSGA-II + UCB1   & 1.86 $\pm$ 0.71 & 0.17 $\pm$ 1.31    & 181.4 $\pm$ 15.9 & 54.4 $\pm$ 7.7 & 0.378 & 0.038 & 0.064 & 13.7 $\pm$ 2.8 \\
 & D-UCB1 myopic      & 1.81 $\pm$ 0.75 & 0.16 $\pm$ 1.26    & 181.4 $\pm$ 15.8 & 54.5 $\pm$ 7.9 & 0.379 & 0.039 & 0.061 & 13.6 $\pm$ 3.2 \\
 & D-LinUCB myopic    & 1.83 $\pm$ 0.71 & 0.15 $\pm$ 1.33    & 181.5 $\pm$ 15.9 & 54.2 $\pm$ 7.5 & 0.378 & 0.039 & 0.063 & 13.8 $\pm$ 3.0 \\
 & D-UCB1 batched     & 8.36 $\pm$ 1.02 & 24.39 $\pm$ 14.58  & 138.0 $\pm$ 30.9 & 41.6 $\pm$ 11.0 & 0.349 & 0.033 & 0.057 & 18.1 $\pm$ 4.4 \\
 & Myopic ILP         & 0.00            & 0.00               & 181.7 $\pm$ 15.7 & 53.9 $\pm$ 7.7 & 0.390 & 0.051 & 0.092 & 14.4 $\pm$ 3.3 \\
 & Lexicographic rule & 2.12 $\pm$ 1.38 & 0.51 $\pm$ 1.55    & 180.8 $\pm$ 15.5 & 65.1 $\pm$ 8.0 & 0.390 & 0.087 & 0.091 & 7.5 $\pm$ 3.0 \\
 & Two-stage (sum)    & 0.00            & $-$0.37 $\pm$ 1.21 & 182.4 $\pm$ 15.9 & 55.7 $\pm$ 7.7 & 0.377 & 0.038 & 0.054 & 12.8 $\pm$ 3.3 \\
 & Two-stage ($F_3$)  & 0.00            & 0.08 $\pm$ 1.13    & 181.6 $\pm$ 15.7 & 56.0 $\pm$ 7.7 & 0.390 & 0.037 & 0.060 & 13.0 $\pm$ 3.2 \\
\midrule
18 & S-NSGA-II + UCB1   & 1.91 $\pm$ 0.70 & 0.31 $\pm$ 1.31    & 183.3 $\pm$ 15.8 & 56.1 $\pm$ 8.0 & 0.388 & 0.034 & 0.056 & 7.1 $\pm$ 2.4 \\
 & D-UCB1 myopic      & 1.91 $\pm$ 0.71 & 0.27 $\pm$ 1.30    & 183.4 $\pm$ 15.8 & 56.0 $\pm$ 8.0 & 0.389 & 0.034 & 0.054 & 7.1 $\pm$ 2.5 \\
 & D-LinUCB myopic    & 1.90 $\pm$ 0.75 & 0.41 $\pm$ 1.34    & 183.2 $\pm$ 15.8 & 55.8 $\pm$ 7.7 & 0.389 & 0.034 & 0.055 & 7.2 $\pm$ 2.4 \\
 & D-UCB1 batched     & 8.10 $\pm$ 0.94 & 12.77 $\pm$ 7.73   & 160.4 $\pm$ 19.0 & 49.1 $\pm$ 7.7 & 0.389 & 0.033 & 0.053 & 8.1 $\pm$ 2.6 \\
 & Myopic ILP         & 0.00            & 0.00               & 183.9 $\pm$ 15.8 & 55.6 $\pm$ 8.0 & 0.400 & 0.045 & 0.080 & 7.9 $\pm$ 2.7 \\
 & Lexicographic rule & 2.41 $\pm$ 1.66 & 0.66 $\pm$ 1.70    & 182.7 $\pm$ 15.6 & 67.0 $\pm$ 8.2 & 0.399 & 0.088 & 0.079 & 3.1 $\pm$ 2.0 \\
 & Two-stage (sum)    & 0.00            & $-$0.28 $\pm$ 1.23 & 184.4 $\pm$ 15.8 & 57.1 $\pm$ 7.8 & 0.386 & 0.033 & 0.048 & 6.8 $\pm$ 2.6 \\
 & Two-stage ($F_3$)  & 0.00            & 0.00 $\pm$ 1.30    & 183.9 $\pm$ 15.6 & 57.0 $\pm$ 7.6 & 0.399 & 0.032 & 0.050 & 7.1 $\pm$ 2.5 \\
\bottomrule
\end{tabular}
\end{table*}

\subsection{RQ1: Evolution of the Price of Fairness}
\label{sec:results:rq1}

The top row of Fig.~\ref{fig:clearing} shows $\pof$ round by round (solid lines). With myopic clearing,
the framework's mean $\pof$ was 1.60\%, 1.81\% and 1.91\% for
$\tau = 6$, 12 and 18. The first round, which clears the initial
backlog of 75 vertices, cost 8.44\%; after that the mean fell to
1.28--1.61\%. Over all rounds the per-run slope was negative
($p < 0.001$), but only because of the first round: over rounds 2--25
it did not differ from zero ($p = 0.42$--0.54). Once the backlog is
cleared, $\pof$ is stable over time.

What drives $\pof$ is the size of each round, not the passage of
time. Under the 90\% floor, $\pof$ can be positive only when the ILP
optimum reaches 10 transplants, which happened in 21--24\% of rounds.
In those rounds the framework gave up 7.7--7.8\% of transplants; in
all others it implemented a maximum matching.

The per-round price is also largely temporary. Summed over
the 25 rounds, the framework gave up 7.0, 7.7 and 8.0 transplants per
run against each round's optimum, yet it finished only $-0.06$, 0.31
and 0.50 transplants behind the myopic ILP. In other words,
94--100\% of the shortfall against each round's optimum was recovered
in later rounds, and $\Delta$ against the myopic ILP was not
significant at any $\tau$: patients who are not matched in one round
usually remain in the pool for the next. This holds within the synthetic
model; it does not establish that the same would occur in the ANZKX
program.

\subsection{RQ2: Myopic versus Batched Clearing}
\label{sec:results:rq2}

Batched clearing with $n^* = 90$ was worse at every $\tau$
(Fig.~\ref{fig:clearing}). At $\tau = 6$ the pool never reached 90
patients, so no one was transplanted and 65.6 highly sensitized
patients per run withdrew, against 25.5 under myopic clearing. At
$\tau = 12$ it transplanted 24\% fewer patients (138.0 vs.\ 181.4,
$d_z = 1.72$) and 24\% fewer highly sensitized patients, and 33\% more
of them withdrew. At $\tau = 18$ both losses were 12\% (all
$p_{\mathrm{bonf}} < 0.001$). Each batch was large, so $\pof$ sat near
the floor in every cleared round (8.1--8.4\%), while $F_3$ and $F_4$
were similar to myopic clearing. Waiting did not improve fairness and reduced transplants.

Better-tuned batching narrowed the gap without closing it. Forcing a clearing at least every three
rounds avoided the failure at $\tau = 6$ and cut the loss to about
4.5\% at $\tau = 12$ and 18. Lowering the threshold to $n^* = 60$
reduced it to 5.5\% and 1.5\%. No variant transplanted more patients
than myopic clearing. This agrees with the near-optimality of greedy
matching without departure information~\cite{anderson2017efficient,akbarpour2020thickness}, and, in this model, extends to the fairness
objectives.

\begin{figure*}[!htbp]
\centering
\includegraphics[width=\textwidth]{fig4_myopic_vs_batched.png}
\caption{Myopic (solid) versus batched (dashed) clearing: per-round
$\pof$ (top) and cumulative transplants (bottom), mean of 100 runs.}
\label{fig:clearing}
\end{figure*}

\subsection{RQ3: Comparison with the Baselines}
\label{sec:results:rq3}

\emph{Myopic ILP.} The framework transplanted the same number of
patients at every $\tau$ ($p_{\mathrm{bonf}} = 1.00$) and was fairer
on all three fairness objectives ($p_{\mathrm{bonf}} < 0.001$;
Fig.~\ref{fig:forest}). It reduced $F_3$ by 21--24\%
($|d_z| = 1.4$--1.7), $F_4$ by 28--34\% ($|d_z| = 1.3$--1.8) and
$F_2$ by 2.7--3.9\% ($|d_z| = 1.9$--2.7), and 3--10\% fewer highly
sensitized patients withdrew ($p_{\mathrm{bonf}} \leq 0.032$). Part of this gain needs no loss of volume: with a 100\% floor ($\pof = 0$ in every
round) the framework still matched the ILP's volume ($\tau = 12$, 30 runs, Section~\ref{sec:results:rq4}: 179.2 vs.\ 178.4) and reduced $F_3$ by 18\% ($p_{\mathrm{bonf}} < 0.001$). The
volume ILP leaves fairness unused among the maximum matchings; the
two-stage ILP below exploits it directly.

\emph{Two-stage ILP.} The two-stage ILP found fairer maximum-volume
matchings than the framework at all three levels of comparison
(Table~\ref{tab:levels}). On the 75-vertex identical pools 67\% of its
stage-2 solves reached the time limit, so these results are best-found
solutions; on trajectory rounds the pools are smaller and only 4.0\% of
solves reached the limit (Section~\ref{sec:results:cost}).
\begin{itemize}
  \item \emph{Identical pools.} The two-stage ILP (sum) reached
    $F_3 = 0.0217$ and $F_4 = 0.0348$, against 0.0548 and 0.1143 for
    the framework's maximum-volume choice (both
    $p_{\mathrm{bonf}} \leq 0.006$). The framework's front held only
    1.1 solutions at maximum volume per pool on average, and its
    lowest $F_3$ at maximum volume equalled the best $F_3$ found by the
    two-stage ILP ($F_3$) in half of the pools.
  \item \emph{Fixed trajectories.} On the same pools in every round
    ($\tau = 12$; 689 of 750 rounds had a feasible matching), the gap was smaller. On the stage-2
    score~\eqref{eq:twostage}, the volume ILP scored 2.753, the
    two-stage ILP 2.356 and the framework's maximum-volume choice
    2.681. The framework's lowest $F_3$ at maximum volume equalled the
    best $F_3$ found by the two-stage ILP ($F_3$) in 95.4\% of rounds
    (100\% of enumerated rounds), so most of the remaining gap comes from
    the selection rule, which differs from the stage-2 objective, rather
    than from the search.
  \item \emph{Full trajectories.} Over 100 paired runs, against the
    framework with its 90\% floor, the two-stage ILP made 1.3, 1.0 and
    1.0 more transplants per run ($\tau = 6$, 12, 18; all
    $p_{\mathrm{bonf}} < 0.001$) and 1.1--1.7 more highly sensitized
    transplants ($p_{\mathrm{bonf}} < 0.001$). Its $F_4$ was lower at
    every $\tau$ ($p_{\mathrm{bonf}} < 0.001$), its $F_2$ was lower at
    $\tau = 12$ and 18 ($p_{\mathrm{bonf}} \leq 0.040$), and its $F_3$ was
    lower at $\tau = 6$ ($p_{\mathrm{bonf}} = 0.034$) and did not differ
    at $\tau = 12$ and 18. With a 100\% floor, the framework matched the
    two-stage ILP's volume but had higher $F_2$, $F_3$ and $F_4$
    ($\tau = 12$, 30 runs, $p_{\mathrm{bonf}} \leq 0.004$).
\end{itemize}
Within the time limits used here, the two-stage ILP is therefore the
stronger approach when the aim is the fairest matching at maximum volume. Against the myopic ILP, the two-stage ILP was much fairer
and made as many or more transplants (more at $\tau = 6$,
$p_{\mathrm{bonf}} < 0.001$).

\emph{Tie-breaking.} In 956 enumerated rounds, 24\% had more than one
optimal maximum matching, but the tied matchings differed in $F_2$, $F_3$
or $F_4$ in only 0.1\% of rounds, so the HLA tie-break rarely affects the
results.

\emph{Lexicographic rule}~\cite{dickerson2014price}. Volume and $\pof$ did
not differ from the framework. The lexicographic rule transplanted about
20\% more highly sensitized patients and lost fewer of them to withdrawal
(3.1--17.5 vs.\ 7.1--25.5 per run), but their match rate rose well above
everyone else's, giving the largest $F_3$ of all myopic methods
(0.087--0.095). The two methods pursue different goals, maximum priority
versus balanced access, and the choice between them is a policy decision.

\emph{Static baselines.} Static NSGA-II with UCB1 and static NSGA-III with
a fixed rate gave the same volume, $\pof$ and fairness as the dynamic
framework; warm-starting added only 0.55\% hypervolume at $\tau = 18$
($p_{\mathrm{bonf}} = 0.007$).

\begin{table}[!htbp]
\caption{Framework and ILP baselines at three levels (means, $\tau = 12$
for B and C). A: identical pools; B: fixed trajectories; C: full
trajectories. Score: stage-2 objective~\eqref{eq:twostage}.}
\label{tab:levels}
\begin{tabular}{lcccc}
\toprule
Method & $F_2$ & $F_3$ & $F_4$ & Score \\
\midrule
\GROUP{A. Identical pools} \\
Volume ILP          & 0.350 & 0.056 & 0.126 & 2.875 \\
Two-stage ILP (sum) & 0.324 & 0.022 & 0.035 & 1.826 \\
Two-stage ILP ($F_3$) & 0.347 & 0.022 & 0.047 & 1.977 \\
Framework, 100\% pick & 0.348 & 0.055 & 0.114 & 2.783 \\
Framework, 90\% pick  & 0.330 & 0.023 & 0.055 & 2.212 \\
\GROUP{B. Fixed trajectories} \\
Volume ILP          & 0.382 & 0.049 & 0.090 & 2.753 \\
Two-stage ILP (sum) & 0.378 & 0.036 & 0.054 & 2.356 \\
Two-stage ILP ($F_3$) & 0.382 & 0.035 & 0.058 & 2.414 \\
Framework, 100\% pick & 0.379 & 0.041 & 0.069 & 2.681 \\
Framework, 90\% pick  & 0.378 & 0.037 & 0.061 & 2.578 \\
\GROUP{C. Full trajectories} \\
Myopic ILP          & 0.390 & 0.051 & 0.092 & -- \\
Two-stage ILP (sum) & 0.377 & 0.038 & 0.054 & -- \\
Two-stage ILP ($F_3$) & 0.390 & 0.037 & 0.060 & -- \\
Framework, 100\% floor & 0.381 & 0.041 & 0.069 & -- \\
Framework, 90\% floor  & 0.379 & 0.039 & 0.061 & -- \\
\bottomrule
\end{tabular}
\NOTE{Score in B: mean of per-run means over all rounds.}
\NOTE{C: 100 paired runs, except the 100\% floor (30 runs).}
\end{table}

\begin{figure*}[!htbp]
\centering
\includegraphics[width=\textwidth]{fig5_forest_tau12.png}
\caption{Framework minus each baseline at $\tau = 12$ (mean and 95\% CI,
100 runs). Negative values favor the framework; filled markers are
significant after Bonferroni correction.}
\label{fig:forest}
\end{figure*}

\subsection{RQ4: Trade-off Solutions below Maximum Volume}
\label{sec:results:rq4}

A single two-stage ILP solution fixes the volume at its maximum. The
front, in contrast, contains matchings at several volumes. We asked
whether these give meaningful options and how close they are to the best
ILP solution found at the same volume (Fig.~\ref{fig:tradeoff}).

\emph{Identical pools.} The front grew from 3.4 solutions at the 95\% level
to 23.0 at 80\%. In 23\%--90\% of pools (95\% to 80\% levels) it held a
solution with a lower stage-2 score than the two-stage ILP at maximum
volume, and in 3\%--60\% a solution at least as good on $F_2$, $F_3$ and
$F_4$ together. The two-stage ILP solved at the same volumes still found
fairer solutions at every level (1.291 vs.\ 1.318 at 80\%), but
97\%--100\% of these solves reached the 120~s limit, so they are
best-found rather than proven optimal solutions. Each ILP solve returns
one matching at one volume, whereas the front gave all of them in a
single run of about 27~s.

\emph{Fixed trajectories.} On trajectory rounds, where only 5\%--12\% of
the ILP solves reached the 60~s limit, the front was close to the ILP
reference at every level. At $\tau = 12$ its fairest solution scored
2.416 and 2.167 at the 90\% and 80\% levels, against 2.379 and 2.146, with
the same $F_3$ (0.035 at 90\%; 0.023 vs.\ 0.024 at 80\%). Where the 80\%
level allowed fewer transplants than $F_1^*$ (70\% of rounds), the front
offered 7.5 options on average, and in 64\% of these rounds at least one
was fairer than the two-stage ILP at maximum volume. At this level the
fairest front solution had significantly lower $F_2$, $F_3$ and $F_4$ than
the two-stage ILP at every $\tau$ ($p_{\mathrm{bonf}} < 0.001$). The floor rule did not pick these solutions: over all rounds (Fig.~\ref{fig:tradeoff}B), its choice scored 2.540 at 80\% and
2.830 at maximum volume, against 2.167 and 2.548 for the fairest solution
on the same front.

Three quantities therefore need to be kept apart: the quality of the
front, which was close to the ILP reference; the quality of the matching
that the 90\% floor rule implemented from it, which was clearly worse; and
the quality of the two-stage ILP solution at maximum volume. The
framework's contribution lies in the first: it provides trade-off
solutions below maximum volume. The second is a limitation of the current
selection rule.

\emph{Selection rules over full trajectories.} The rules move the
framework along a consistent trade-off at all three $\tau$
(Table~\ref{tab:select}). The 80\% floor gave a
lower $F_3$ than the two-stage ILP at every $\tau$ (0.045, 0.028 and
0.025 against 0.054, 0.035 and 0.032; $p_{\mathrm{bonf}} < 0.001$) and
a significantly lower $F_2$ and $F_4$ at $\tau = 6$ and 12, for 2.1--2.9\% fewer
transplants ($\Delta = 1.7$--2.1\%). The compromise rule reduced $F_3$
to 0.013--0.026 at a much larger cost ($\pof = 36$--37\%,
$\Delta = 5.7$--6.6\%). Both rules also transplanted fewer highly
sensitized patients than the two-stage ILP (2.2--2.6 fewer with the
80\% floor), because $F_3$ can fall when fewer patients are matched
overall. These options are not available from a single two-stage ILP
solution; whether a given trade-off is acceptable is a policy choice.

\begin{figure*}[!htbp]
\centering
\includegraphics[width=\textwidth]{fig6_tradeoff_levels.png}
\caption{Fairness against volume level $\rho$: best front solution,
floor-rule choice, two-stage ILP (best found, time-limited) and
two-stage ILP at maximum volume (dotted). Mean and 95\% CI. Panel B averages all rounds with $F_1^{*} > 0$ of the 30 runs; part~B of Table~\ref{tab:levels} averages per-run means, so the two scores are not directly comparable.}
\label{fig:tradeoff}
\end{figure*}

\begin{table*}[!htbp]
\caption{Selection rules over full trajectories compared with the
two-stage ILP and the myopic ILP (mean over the same 30 paired runs).}
\label{tab:select}
\begin{tabular}{lccccccc}
\toprule
Configuration & Tx & High-cPRA tx & $F_2$ & $F_3$ & $F_4$ & PoF (\%) & $\Delta$ (\%) \\
\midrule
\GROUP{$\tau = 6$} \\
Framework, 100\% floor & 171.7 & 50.6 & 0.359 & 0.065 & 0.106 & 0.00 & $-$1.02 \\
Framework, 90\% floor  & 170.2 & 49.9 & 0.355 & 0.060 & 0.101 & 1.46 & $-$0.13 \\
Framework, 80\% floor  & 166.9 & 49.3 & 0.345 & 0.045 & 0.077 & 9.86 & 1.81 \\
Framework, compromise  & 158.9 & 46.6 & 0.320 & 0.026 & 0.047 & 36.04 & 6.55 \\
Two-stage ILP (sum)    & 171.9 & 51.5 & 0.355 & 0.054 & 0.091 & 0.00 & $-$1.12 \\
Myopic ILP             & 170.0 & 48.9 & 0.370 & 0.075 & 0.140 & 0.00 & 0.00 \\
\GROUP{$\tau = 12$} \\
Framework, 100\% floor & 179.2 & 54.9 & 0.381 & 0.041 & 0.069 & 0.00 & $-$0.49 \\
Framework, 90\% floor  & 178.6 & 54.9 & 0.378 & 0.037 & 0.061 & 1.68 & $-$0.10 \\
Framework, 80\% floor  & 175.3 & 53.3 & 0.371 & 0.028 & 0.048 & 10.34 & 1.74 \\
Framework, compromise  & 168.3 & 51.1 & 0.357 & 0.015 & 0.028 & 37.09 & 5.70 \\
Two-stage ILP (sum)    & 179.1 & 55.7 & 0.375 & 0.035 & 0.056 & 0.00 & $-$0.43 \\
Myopic ILP             & 178.4 & 53.9 & 0.389 & 0.050 & 0.093 & 0.00 & 0.00 \\
\GROUP{$\tau = 18$} \\
Framework, 100\% floor & 181.8 & 56.7 & 0.392 & 0.037 & 0.063 & 0.00 & $-$0.72 \\
Framework, 90\% floor  & 180.6 & 56.5 & 0.388 & 0.034 & 0.058 & 1.83 & $-$0.05 \\
Framework, 80\% floor  & 176.8 & 54.9 & 0.382 & 0.025 & 0.043 & 10.62 & 2.06 \\
Framework, compromise  & 169.9 & 52.3 & 0.370 & 0.013 & 0.025 & 37.43 & 5.92 \\
Two-stage ILP (sum)    & 181.5 & 57.5 & 0.385 & 0.032 & 0.048 & 0.00 & $-$0.53 \\
Myopic ILP             & 180.5 & 55.5 & 0.400 & 0.045 & 0.081 & 0.00 & 0.00 \\
\bottomrule
\end{tabular}
\NOTE{Tx: transplants.}
\end{table*}

\subsection{RQ5: Operator Selection and ILP Guidance}
\label{sec:results:rq5}
UCB1 and LinUCB were statistically identical on every outcome
($p_{\mathrm{bonf}} = 1.00$). UCB1 spread its choices almost evenly over
the four mutation rates, whereas LinUCB leaned slightly towards
$p_m = 0.25$ under myopic clearing and towards $p_m = 0.01$ under
batched clearing. Table~\ref{tab:ablation} isolates the ILP guidance.
NSGA-III alone never reached the volume optimum, and its selected
solutions had a mean $\pof$ of 28.3\%, against 8.2\% for the full
hybrid ($p_{\mathrm{bonf}} < 0.001$); roughly 20 percentage points of
the unguided $\pof$ were search failure rather than fairness. The seed
alone changed nothing, because the greedy repair rebuilt it into a
smaller matching; repair priority raised optimum attainment to 77\% of
pools, and the archive closed the remaining gap
($p_{\mathrm{bonf}} = 0.007$ for hypervolume). In the trajectory runs (Section~\ref{sec:results:cost}), the volume ILP took 0.52~s per round, 2.7\% of a
framework clearing (19.30~s). Without guidance, no operator setting
reached the optimum; with guidance, every setting did, and
hypervolume differed by less than 3\%. UCB1 was numerically better than fixed rates of 0.10 and 0.25 and slightly below a fixed rate of 0.01. The ILP
guidance, not the operator selection, drives search quality.

\begin{table}[!htbp]
\caption{Ablation of the ILP guidance (UCB1): one round on 30 pools of 75
vertices (ILP optimum 42.3 transplants). PoF of the 90\% floor choice
relative to the ILP optimum.}
\label{tab:ablation}
\begin{tabular}{lcccc}
\toprule
Setting & At optimum & Max $F_1$ & HV & PoF (\%) \\
\midrule
NSGA-III alone  & 0/30  & 32.6 & 0.566 & 28.3 \\
ILP seed only   & 0/30  & 32.7 & 0.566 & 28.8 \\
Seed + priority & 23/30 & 41.8 & 0.696 & 9.2 \\
Full hybrid     & 30/30 & 42.3 & 0.703 & 8.2 \\
\bottomrule
\end{tabular}
\end{table}

\subsection{RQ6: Sensitivity and Scalability}
\label{sec:results:rq6}
Table~\ref{tab:sens} summarizes the sensitivity analyses at
$\tau = 12$. A 12~h CIT limit removed 13--14\% of transplants for every
method, but the framework still matched the ILP and reduced $F_3$.
Changing $\mu_d$ left $\pof$ unchanged; a higher exit rate hurt batched
clearing far more ($-26\%$) than myopic clearing ($-1.6\%$). At
$\lambda = 7$ and 13 the framework again matched the ILP and reduced
$F_3$ by 22\% and 27\%, and $\pof$ rose with pool thickness (1.03\%,
1.67\%, 3.04\%). With initial pools of 200 vertices, the framework took
37~s per round, against 5~s for the ILP, kept the ILP's volume (27.6
vs.\ 27.7 per round) and reduced $F_3$ by 39\%; its $\pof$ rose to 7.0\%
because most rounds exceeded 10 transplants, so the 90\% floor allowed
larger cuts. At 500 vertices the ILP took 38~s per round; the framework
was not run because of memory limits.

\begin{table}[!htbp]
\caption{Sensitivity at $\tau = 12$ (30 paired runs): framework (UCB1,
myopic) against the myopic ILP.}
\label{tab:sens}
\begin{tabular}{lccccc}
\toprule
Variant & PoF (\%) & Tx (fw) & Tx (ILP) & $F_3$ (fw) & $F_3$ (ILP) \\
\midrule
Base                          & 1.67 & 178.6 & 178.4 & 0.037 & 0.050 \\
$\mathrm{CIT}_{\max} = 16$ h  & 1.78 & 178.0 & 177.8 & 0.036 & 0.049 \\
$\mathrm{CIT}_{\max} = 12$ h  & 1.51 & 154.6 & 153.9 & 0.033 & 0.040 \\
$\mu_d = 0.01$                & 1.75 & 179.4 & 179.1 & 0.036 & 0.048 \\
$\mu_d = 0.04$                & 1.69 & 175.7 & 176.0 & 0.040 & 0.054 \\
$\lambda = 7$                 & 1.03 & 137.5 & 137.1 & 0.038 & 0.049 \\
$\lambda = 13$                & 3.04 & 224.7 & 225.2 & 0.036 & 0.049 \\
\bottomrule
\end{tabular}
\NOTE{fw: framework; Tx: transplants.}
\end{table}

\subsection{Computational Cost}
\label{sec:results:cost}
Table~\ref{tab:cost} reports the cost of every method. On trajectory
rounds the two-stage ILP (sum) was faster than the framework (5.08 vs.\
19.30~s per round), and only 4.0\% of its stage-2 solves reached the time
limit. On the 75-vertex identical pools it took 233.3~s against 27.2~s,
with 67\% of solves at the limit, and solves at lower volume levels almost
always reached it. The two-stage ILP is fast once the backlog is cleared,
but its cost grows with pool size and with the number of volume levels,
whereas the framework returns a whole front at a cost that does not depend on the number of volume levels.

\begin{table}[!htbp]
\caption{Computational cost on one core. Pool: one round on a
75-vertex pool (stage-2 limit 300~s). Trajectory: mean seconds per
cleared round at $\tau = 12$ (stage-2 limit 60~s). TL: share of
stage-2 solves that reached the time limit.}
\label{tab:cost}
\begin{tabular}{lcccc}
\toprule
Method & Pool (s) & TL & Trajectory (s) & PrefLib (s) \\
\midrule
Volume (myopic) ILP          & 3.2   & --     & 0.52     & 0.16 \\
Two-stage ILP (sum)          & 233.3 & 0.67   & 5.08     & 2.45 \\
Two-stage ILP ($F_3$)        & 58.5  & 0.17   & 1.60     & 0.76 \\
Framework (incl.\ ILP)       & 27.2  & --     & 19.30    & 4.41 \\
Lexicographic rule           & --    & --     & 1.14     & 0.38 \\
Stage 2, lower levels  & --    & 0.05--0.12$^{a}$ & 24.96$^{b}$ & -- \\
\bottomrule
\end{tabular}
\NOTE{$^{a}$Range over the four lower levels, trajectory rounds.}
\NOTE{$^{b}$All four lower levels per round, mean over $\tau = 6$, 12, 18.}
\end{table}

\subsection{Out-of-Distribution Validation on PrefLib}
\label{sec:results:ood}

Table~\ref{tab:ood} reports 120 paired runs. The framework matched the
volume of the myopic ILP ($p_{\mathrm{bonf}} = 0.19$) and lowered $F_2$
significantly, but its $F_3$ gain was small and not significant. The
two-stage ILP (sum) again found fairer maximum-volume matchings, the
lexicographic rule again over-corrected ($F_3 = 0.080$), and batched
clearing lost 3.9\% of transplants and doubled $F_3$. PrefLib graphs are
dense, so most patients are matched in the first round and few later
rounds offer a choice between maximum matchings. Centers are assigned at
random, so $F_4$ is not reported.

\begin{table}[!htbp]
\caption{PrefLib 00036: 12 instances of 128--147 vertices, 10 paired
runs each, $\tau = 12$, 25 rounds (means). $F_4$ is not reported
because centers are assigned at random.}
\label{tab:ood}
\begin{tabular}{lcccccc}
\toprule
Configuration & Tx & HS tx & $F_2$ & $F_3$ & $\pof$ (\%) & $\Delta$ (\%) \\
\midrule
Framework (UCB1)        & 75.03 & 10.59 & 0.316 & 0.060 & 0.73 & 0.50 \\
Framework (LinUCB)      & 74.96 & 10.62 & 0.315 & 0.059 & 0.77 & 0.62 \\
Framework, batched$^{a}$ & 72.12 & 10.16 & 0.325 & 0.118 & 2.20 & 4.35 \\
Two-stage ILP (sum)     & 75.79 & 11.40 & 0.312 & 0.055 & 0.00 & $-$0.50 \\
Two-stage ILP ($F_3$)   & 75.36 & 11.45 & 0.324 & 0.056 & 0.00 & 0.06 \\
Myopic ILP              & 75.38 & 10.72 & 0.325 & 0.062 & 0.00 & 0.00 \\
Lexicographic rule      & 75.09 & 13.57 & 0.324 & 0.080 & 1.54 & 0.41 \\
\bottomrule
\end{tabular}
\NOTE{Tx: transplants; HS: highly sensitized.}
\NOTE{$^{a}$Three-round maximum interval.}
\end{table}

\subsection{Limitations}
\label{sec:results:lim}
All results come from synthetic simulation informed by public ANZKX
information, not from patient-level clinical data. The exit rate,
withdrawal thresholds, transport and CIT constraints and the HLA model
are modeling assumptions rather than validated ANZKX parameters, and no
single arrival rate matches both the ANZKX volume and pool size,
although the findings held at $\lambda = 7$, 10 and 13. $F_4$ is a
proxy. Both approaches have limits. The two-stage ILP formulation is exact, but its
solves were time-limited: on 75-vertex pools most stopped at the limit
(67\% at maximum volume, 97\%--100\% at lower volumes), so those
solutions are the best found, not proven optima, and each solve returns
one matching at one volume. The framework costs about 19~s per round,
was not run at 500 vertices because of memory limits, and its front is
approximate. $\pof$ is bounded by the 90\% floor, and the floor rule,
which differs from the stage-2 objective, did not always choose the
fairest solution on the front. Batched clearing was tested without
information about upcoming departures, and the framework's fairness
gain over the volume ILP was small on PrefLib.

\section{Conclusion}
\label{sec:conclusion}
This paper asked how the price of fairness behaves when a kidney
exchange is cleared month after month rather than once, whether
myopic or batched clearing better preserves the balance between
efficiency and equity, and what a multi-objective search adds to a
two-stage integer programming benchmark. We studied this in an
ANZKX-informed simulator with withdrawals, deceased-donor exits and cold
ischemia limits across six regional nodes, using an ILP-guided NSGA-III
with bandit-based operator selection.

The results lead to four findings. First, the per-round price of
fairness was small (1.6--1.9\%) and did not accumulate over the
simulated trajectory, while the access, logistics and waiting-time gaps
all narrowed without a significant loss of total volume. Second, batching lost
transplants without improving fairness. Third, the two-stage ILP, solved
within practical time limits, found fairer maximum-volume matchings than
the framework, although on the larger pools these were best-found rather
than proven optimal solutions. Fourth, below maximum volume the front contained trade-offs close to the ILP reference on trajectory rounds, but the 90\% floor rule did not always implement the fairest of them; ILP guidance, not the bandit, drove the quality of the search.

Within this computational model, matching every month outperformed
waiting for a thicker pool. When the aim is the fairest matching at maximum
volume, the two-stage ILP is the better tool, provided its solves can run
to completion or close to it; when some volume may be given up for
fairness, the front provides a range of options from which a policy can
be chosen. These are model results. They do not establish that the same
outcomes would occur in ANZKX, and they require validation with program
data before any practical or clinical conclusion can be drawn. Future
work should test the methods on registry data, subject to ethics
approval, design selection rules aligned with the stage-2 objective,
compare the front with exact $\varepsilon$-constraint fronts, and study
clearing policies that anticipate departures.

\REFERENCES

\end{document}
